\section{Why this matters for measurement}
\label{sec:measurement}

In educational measurement terms, this is a method for writing test items and scoring them automatically, at the same time.

\begin{figure*}[t]
\centering
\begin{tikzpicture}[
    font=\small,
    tmpl/.style={draw, rounded corners, fill=blue!8, align=center,
                 text width=72mm, minimum height=11mm, inner sep=4pt},
    ar/.style={-{Stealth[length=2mm]}, thick, gray!55!black},
]
% --- the item model ---
\node[tmpl] (t) at (4.4,3.05) {\textbf{One item model}\\[1pt]
  {\scriptsize\ttfamily triangle\_with\_altitude(angle\_A${=}\alpha$,\ angle\_B${=}70^\circ$)}\\[2pt]
  {\scriptsize vary one parameter: $\alpha \in \{30^\circ,\,40^\circ,\,50^\circ\}$}};
% --- three generated instances ---
\foreach \sx/\cx/\cy/\ang in {0/1.653/0.955/30, 3.4/1.532/1.286/40, 6.8/1.395/1.662/50}{
  \begin{scope}[shift={(\sx,0)}]
    \coordinate (A) at (0,0); \coordinate (B) at (2,0);
    \coordinate (C) at (\cx,\cy); \coordinate (H) at (\cx,0);
    \draw[thick] (A)--(B)--(C)--cycle;
    \draw[densely dashed] (C)--(H);
    \draw (\cx-0.16,0)--(\cx-0.16,0.16)--(\cx,0.16); % right-angle mark at H
    \draw (0.45,0) arc (0:\ang:0.45);                % angle arc at A
    \node[font=\scriptsize] at (\ang/2:0.74) {$\ang^\circ$};
    \foreach \q in {A,B,C,H}{\fill (\q) circle (1.1pt);}
    \node[below left=-3pt, font=\scriptsize]  at (A) {$A$};
    \node[font=\scriptsize, anchor=north west] at ([yshift=-1pt]B) {$B$};
    \node[above=-2pt, font=\scriptsize]       at (C) {$C$};
    \node[font=\scriptsize, anchor=north east] at ([yshift=-1pt]H) {$H$};
    \node[font=\small] at (1,-0.62)
      {$\alpha=\ang^\circ$\ \ {\color{green!45!black}\checkmark}};
  \end{scope}
}
% --- fan-out arrows + verdict line ---
\foreach \sx in {0,3.4,6.8}{
  \draw[ar] (t.south) .. controls +(0,-0.5) and +(0,0.8) .. (\sx+1,1.55);}
\node[font=\small\itshape, color=gray!40!black] at (4.4,-1.18)
  {all three verified automatically, none by a human};
\end{tikzpicture}
\caption{From one item model to a family of verified items. This is automatic item generation with the answer key built in. Hold every parameter fixed but the angle at $A$ and sweep it through $30^\circ, 40^\circ, 50^\circ$. That produces parallel variants. For each, the checker of \cref{sec:scoring} confirms that the altitude from $C$ meets $AB$ at a right angle (\checkmark), with no human review. One model yields a refreshable supply of items, each carrying its own key.}
\label{fig:variants}
\end{figure*}

Each template is an item generator. Fixing every parameter but one and sweeping that parameter, for example an angle through $30^\circ, 40^\circ, 50^\circ$, produces a family of structurally parallel questions, and the checker confirms each variant without human review (\cref{fig:variants}). The answer key is produced \emph{with} each item rather than written afterward. So a common source of scoring error, a hand-written key that has drifted from its question, is largely designed out. In the vocabulary of the field the template is an \emph{item model} \citep{gierl2012itemmodels}, now one that emits its own key. That move also shifts where quality control happens, in the way the item-generation literature recommends. Rather than vetting every generated item, one reviews the construction logic that produces them all \citep{gierl2016review}.

The diagram a student (or a model) produces is a constructed response. The same description that built the item can score it. Here a measurement audience should consider two claims separately. The \emph{generation} claim rests on firmer ground. The key is computed with the item rather than written separately, so question and key are far less likely to drift apart. And the key is a short, readable list of exactly what must be true (\cref{tab:predicates}), open to inspection rather than embedded in a holistic rubric. The \emph{scoring} claim must be established empirically, and our validity study (\cref{sec:scoring}) shows that it is not yet established. The auto-score agrees with human graders only moderately and errs in a single direction. For now it is an optimistic estimate. That is fit for benchmarking and for flagging likely-wrong work. It is not yet fit for standalone high-stakes scoring. Closing that gap means adding the numeric checks that \cref{sec:scoring} pinpoints.

\paragraph{Three measurement questions the checklist opens.}
First, isomorphism. Variants minted from one item model are parallel in construction. Parallel construction does not guarantee equal difficulty. A $30^\circ$ angle may be harder to draw than a $50^\circ$ one. Whether the variants are true isomorphs is a calibration question. The framework makes that study inexpensive, because every variant comes with its key. Second, partial credit. The checker decides every requirement separately, so a score need not be all-or-nothing. A diagram that places the altitude correctly but misses the stated angle can earn credit for what it got right. The per-requirement verdicts provide the basis for a partial-credit rubric. The weights and score categories remain to be decided. Our graders also rated diagrams requirement by requirement. They agreed with the checker on 98\% of 864 individual ratings. The overall disagreement in \cref{sec:scoring} comes largely from requirements the checklist left out, not from the ones it contains. Third, the checklist as a Q-matrix. Each item's list of checks names the geometric requirements the item imposes, in a fixed vocabulary shared across items. That is the raw form of a Q-matrix, the item-by-attribute table a cognitive diagnostic model needs. Mapping requirements to the skills a student uses still takes judgment and empirical validation. But an initial version of the table is produced by the item model rather than coded by hand afterward.

\paragraph{Scope.}
The version described here covers plane geometry. The same machinery should extend to coordinate graphs and 3D solids by adding new kinds of building block. We have not built those yet. Where a topic has no natural template, an author can still write items and keys by hand and score them with the same checker. The core idea is simple. When a diagram is built from a structured description, that description says what must be checked, and a machine can check it. Whether the description is complete is a question for validation, and \cref{sec:scoring} shows why.
